\documentclass[11pt,a4paper]{article}
\usepackage[T1]{fontenc}
\usepackage[utf8]{inputenc}
\usepackage{lmodern}
\usepackage{amsmath,amssymb}
\usepackage[margin=25mm]{geometry}
\usepackage{longtable,booktabs,array,calc}
\usepackage[hidelinks]{hyperref}
\hypersetup{pdftitle={Gauge crossover: comparison tools and their missing estimates},pdfauthor={navstokgap research working notes}}
\newcommand{\tightlist}{\setlength{\itemsep}{0pt}\setlength{\parskip}{0pt}}
\setlength{\emergencystretch}{3em}
\setcounter{secnumdepth}{0}
\begin{document}
\hypertarget{gauge-crossover-comparison-tools-and-their-missing-estimates}{%
\section{Gauge crossover: comparison tools and their missing
estimates}\label{gauge-crossover-comparison-tools-and-their-missing-estimates}}

\textbf{Maintained conclusion.} The examined comparison and finite-block
techniques provide possible ingredients for crossover control, not an
exhaustive list of ways to prove a physical continuum gap. A normalized,
volume-uniform mixing estimate must be established for the actual
generated interaction. Neither a loop monotonicity statement nor a
finite certification proposal supplies it by itself. Weak-to-strong
transport, cutoff removal and physical reconstruction remain additional
steps.

\hypertarget{correlation-inequalities-what-would-actually-transfer}{%
\subsection{Correlation inequalities: what would actually
transfer}\label{correlation-inequalities-what-would-actually-transfer}}

If Wilson loop expectations increase with \(\beta_W\), an area-law
string tension defined by the corresponding limit is nonincreasing where
the limit and comparison are valid. A lower area-law bound can then
transfer from a weaker to a stronger coupling. Monotonicity of a loop
expectation alone does not imply monotonicity of a connected correlator,
its decay exponent or the vacuum-sector spectral gap. Both terms of a
truncated correlator may increase. The old claim that this closes every
stronger-coupling channel is withdrawn.

Ginibre's method needs the interaction to lie in a cone of products of
single-site functions satisfying the duplicate-variables condition. For
the character cone of \(SU(2)\) the single-site condition reduces to
\[\int\!\!\int\prod_{i=1}^n\big(\chi_i(x)-\chi_i(y)\big)\,dx\,dy
=\sum_{S\subseteq\{1..n\}}(-1)^{|S^c|}\,N_S\,N_{S^c},\qquad
N_S=\dim\operatorname{Hom}\Big(1,\bigotimes_{i\in S}V_i\Big),\] which
vanishes for odd \(n\), equals
\(2N_{1234}+2(\delta_{12}\delta_{34}+\delta_{13}\delta_{24}+\delta_{14}\delta_{23})\ge0\)
for \(n=4\), and is positive in every case computed at \(n=6\) (six spin
\(\tfrac12\): \(70\); six spin \(1\): \(100\); four spin \(\tfrac12\)
and two spin \(1\): \(28\); six spin \(2\): \(260\)). The single-site
condition is therefore no obstruction as far as checked. The obstruction
is that the plaquette term
\(\operatorname{Re}\operatorname{tr}(U_1U_2U_3U_4)\) is not a
nonnegative combination of products of single-link functions: the
abelian proof writes \(\cos(\theta_1+\dots)\) through sums and
differences of angles, which uses the commutativity of the group, and no
analogue exists for a matrix product. The second Griffiths inequality
was not established by the sources read here for \(SU(2)\) and \(SU(3)\)
lattice gauge theory, ; no proof or counterexample is supplied here, and
the centre-based inequalities that do exist are listed in
\href{reasons-to-stop-as-research.md\#tools-and-source-scope}{the
what-would-unblock note} §1.

The finite checks are not an all-order character-cone theorem. The
obstacle identified here is the absence of the needed positive plaquette
decomposition in the examined construction, not a proof that every
correlation inequality for a nonabelian theory is impossible.

\hypertarget{tools-and-source-scope}{%
\subsection{Tools and source scope}\label{tools-and-source-scope}}

These are the sources previously cited in this family, at
\textbf{metadata level} unless a linked canonical proof records a deeper
read. They are leads and scope comparisons, not independently verified
universal impossibility claims.

\begin{longtable}[]{@{}
  >{\raggedright\arraybackslash}p{(\columnwidth - 4\tabcolsep) * \real{0.3333}}
  >{\raggedright\arraybackslash}p{(\columnwidth - 4\tabcolsep) * \real{0.3333}}
  >{\raggedright\arraybackslash}p{(\columnwidth - 4\tabcolsep) * \real{0.3333}}@{}}
\toprule\noalign{}
\begin{minipage}[b]{\linewidth}\raggedright
Tool/source
\end{minipage} & \begin{minipage}[b]{\linewidth}\raggedright
Potential input
\end{minipage} & \begin{minipage}[b]{\linewidth}\raggedright
Missing transfer
\end{minipage} \\
\midrule\noalign{}
\endhead
\bottomrule\noalign{}
\endlastfoot
Ginibre, CMP 16 (1970) 310; Aizenman--Barsky--Fernández, JSP 47 (1987)
343 & Positive cones and Ising-type sharpness & A nonabelian plaquette
interaction satisfying the required hypotheses \\
Mack--Petkova, Ann. Phys. 123 (1979) 442 and 125 (1980) 117;
Tomboulis--Yaffe, CMP 100 (1985) 313 & Centre, electric-flux and
vortex/loop comparisons & A vacuum-sector physical spectral bound; a
string tension is a different observable \\
Fröhlich--Israel--Lieb--Simon, CMP 62 (1978) 1 & Reflection-positive
chessboard bounds for specified events & Good-region mixing and a
normalized block free-energy estimate \\
Fröhlich--Simon--Spencer, CMP 50 (1976) 79 &
Gaussian-domination/infrared techniques for their stated spin
interactions & The needed domination for the nonquadratic plaquette
product; nonabelian target curvature alone is not the failure reason \\
Knabe, JSP 52 (1988) 627 & A finite-size gap criterion under
frustration-free hypotheses & Those hypotheses are not verified for
interacting KS at general coupling; this does not exhaust all
finite-size criteria \\
\href{intermediate-region-finite-verification.md}{Finite mixing
proposal}, \href{dobrushin-uniqueness-wilson.md}{Dobrushin proof} & A
certified supremum over boundary conditions could imply a
Wilson-transfer gap & No certificate at the physical crossover or for
the generated interaction is supplied \\
\href{flow-conjugation-truncation.md}{Exact flow conjugation} &
Spectrum-preserving change of variables & New effective decomposition,
relative error and iterative control \\
\href{weak-coupling-feshbach-reduction.md}{Small-volume bridge} &
Compressed zero-mode estimates with a named complement/Schur target &
Renormalized cutoff-uniform transfer and passage to large physical
volume \\
\end{longtable}

A signed spin-foam representation does not make the original Wilson
probability weight negative, and positive plaquette coefficients do not
make every recoupling amplitude positive. Centre comparisons,
electric-flux sectors and glueball correlations must name the sector
they control. The 2007 approximate centre-decimation proposal
(arXiv:0707.2179, metadata) was not established as a theorem by this
repository's reading.

Gauge fixing can change a conditional specification; merely subtracting
a maximal-tree fraction from a neighbour count does not prove a new
Dobrushin coefficient. The earlier numerical suggestion
\(444\mapsto330\) is unsupported here. Likewise, neither a perturbative
four-dimensional Langevin construction nor its obstruction is proved by
this source survey; an auxiliary stochastic clock is not a physical
energy gap.

\hypertarget{finite-certification-and-the-actual-interaction}{%
\subsection{Finite certification and the actual
interaction}\label{finite-certification-and-the-actual-interaction}}

The earlier illustrative boxes at \(\beta_W\simeq6\) had roughly
\(216\)--\(2000\) links, hence thousands of integration variables. These
are scale estimates from the \href{confinement-scale-bands.md}{bands
note}, not a complexity theorem excluding every certification method.
Monte Carlo by itself is not a certified supremum over boundary
conditions. A rigorous polymer enumeration needs an analytic tail and
its own convergence region. The explicit Wilson criteria here prove
sufficient strong-coupling regions, not limits on all possible future
methods. The \href{wilson-strong-coupling-explicit.md}{strong-coupling
proofs} and \href{dobrushin-uniqueness-wilson.md}{Dobrushin criterion}
own their constants.

The weak side has explicit but restricted one-step tests in
\href{small-field-step-gaussian.md}{Gaussian fluctuation control} and
\href{small-field-step-decay-and-threshold.md}{decay/remainder
thresholds}. The \href{su2-midplane-small-field.md}{mid-plane notebook}
later proves different retained-barrier weak-strip estimates. Their
coefficients and cutoff/volume hypotheses cannot be silently substituted
for a physical-coupling full integral.

After blocking, all generated interactions must be controlled. There is
no proved arrival in the completely analytic set at order-one coupling
and no proved universal statement that all difficulty is one fixed
unknown lattice interaction. A possible mixing certificate at that scale
would still need stable transport from the weak trajectory, physical
spectral visibility, nontriviality and continuum reconstruction.

\hypertarget{consequence-for-state}{%
\subsection{Consequence for STATE}\label{consequence-for-state}}

The productive questions are the exact normalized estimates still
missing, listed in the \href{halving-atlas.md}{gauge atlas}. The older
claim of exactly two unblocking routes, the unsupported closed
large-field tail and flow sharpness, and the assertion that all other
methods fail are withdrawn. The retained source evidence explains what
each examined method would require; it changes no open target into a
theorem and initiates no new research during consolidation.
\end{document}
